%%%%%%%%%%%%%%%%
%%% lev 18 %%%
%%%%%%%%%%%%%%%%

\Chapter{18}

\Verse{1}
{
	\hP{וַיְדַבֵּר יְהֹוָה אֶל מֹשֶׁה לֵּאמֹר׃}
}
{
	\eP{18 And YHWH spoke to Moses, saying,}
}
{
}

\Verse{2}
{
	\hP{דַּבֵּר אֶל בְּנֵי יִשְׂרָאֵל וְאָמַרְתָּ אֲלֵהֶם אֲנִי יְהֹוָה אֱלֹהֵיכֶם׃}
}
{
	\eP{
		“Speak to the children of Israel, and you shall say to
		them: I am YHWH, your God.
	}
}
{
}

\Verse{3}
{
	\hP{כְּמַעֲשֵׂה אֶרֶץ מִצְרַיִם אֲשֶׁר יְשַׁבְתֶּם בָּהּ לֹא תַעֲשׂוּ וּכְמַעֲשֵׂה אֶרֶץ כְּנַעַן אֲשֶׁר אֲנִי מֵבִיא אֶתְכֶם שָׁמָּה לֹא תַעֲשׂוּ וּבְחֻקֹּתֵיהֶם לֹא תֵלֵכוּ׃}
}
{
	\eP{
		You shall shall not do like what is done in the land of
		Egypt, in which you lived; and you shall not do like what
		is done in the land of Canaan, to which I’m bringing you;
		and you shall not go by their laws.
	}
}
{
}

\Verse{4}
{
	\hP{אֶת מִשְׁפָּטַי תַּעֲשׂוּ וְאֶת חֻקֹּתַי תִּשְׁמְרוּ לָלֶכֶת בָּהֶם אֲנִי יְהֹוָה אֱלֹהֵיכֶם׃}
}
{
	\eP{
		You shall do my judgments, and you shall observe my laws,
		to go by them. I am YHWH, your God.
	}
}
{
%	\footnote{
%		You shall do my judgments, and you shall observe my laws. Why is this
%		said now, after we have already had seventeen chapters of laws in
%		Leviticus? This juncture at the beginning of chapter 18 is a turning
%		point: Until now the laws of Leviticus have had to do with ritual:
%		sacrifices, permitted and forbidden food, leprosy. But the rest of
%		Leviticus (especially chapter 19) will give ethical laws as well—not
%		separately, but interspersed with ritual laws. Turning to the ethical
%		laws, we find that, on one hand, they are easier to deal with than
%		ritual laws, first, simply because they are fewer and, second, because
%		the reasons behind ethical law tend to be more readily apparent than
%		these of ritual. We are inclined to question why cows are permitted for
%		food and pigs are not, but when we read a command that “You shall not
%		steal, and you shall not lie, and you shall not act false” (19:11), we
%		do not generally feel the need to ask, “Why not?” On the other hand,
%		what is harder about dealing with ethical laws is that they are not as
%		easily grouped by types, and they in fact usually do not occur in long
%		sections of related laws as do the sacrificial, food, purity, or leprosy
%		laws. The exceptions are the lists of forbidden sexual relations
%		(18:6–30; 20:10–20) and a group of economic laws (Leviticus 25). The
%		ethical laws include such diverse matters as leaving portions of
%		harvests for the poor, judging cases without bias (toward the rich or
%		toward the poor), honoring the aged, fearing one’s parents, not injuring
%		the blind or deaf, not gossiping, not making one’s daughter a
%		prostitute, and not committing adultery. If there is any common
%		denominator to these diverse laws it is: feeling empathy with other
%		human beings and attributing value to the lives of all of them.
%	}
}

\Verse{5}
{
	\hP{וּשְׁמַרְתֶּם אֶת חֻקֹּתַי וְאֶת מִשְׁפָּטַי אֲשֶׁר יַעֲשֶׂה אֹתָם הָאָדָם וָחַי בָּהֶם אֲנִי יְהֹוָה׃}
}
{
	\eP{
		And you shall observe my laws and my judgments, which,
		when a human will do them, he’ll live through them! I am
		YHWH.
	}
}
{
%	\footnote{
%		when a human will do them, he’ll live through them! This way of
%		picturing the laws, as a path to life, begins here. It returns as the
%		climax of the Torah in Deuteronomy. The path to the Tree of Life is
%		blocked at the Torah’s beginning, and the way to recover it is
%		emphasized at the Torah’s end. The laws are not presented as a burden
%		but as a blessing. Those who have characterized the law as a weight that
%		no human can possibly bear, as a curse from which one needs to be saved,
%		may insufficiently appreciate the positive, fulfilling quality of the
%		law as experienced by those who live it. It is presented here as
%		something that humans can do. And those who add so much to the law that
%		it does become a weight, forcing those who observe it into separation
%		from fellow Jews and from the world, risk diminishing the positive force
%		of the law and its value as being ultimately a source of blessing for
%		all the families of the earth.
%	}
}

\Verse{6}
{
	\hP{אִישׁ אִישׁ אֶל כּל שְׁאֵר בְּשָׂרוֹ לֹא תִקְרְבוּ לְגַלּוֹת עֶרְוָה אֲנִי יְהֹוָה׃}
}
{
	\eP{
		“Any man, to any close relative of his: you shall not come
		close to expose nudity. I am YHWH.
	}
}
{
%	\footnote{
%		expose nudity. This means sexual intimacy. In 18:14 the text indicates
%		that, if one has sexual relations with one’s aunt, one exposes one’s
%		uncle’s nudity. So to “expose nudity” must be understood broadly as
%		meaning to enter the zone of another person’s sexuality.
%	}
}

\Verse{7}
{
	\hP{עֶרְוַת אָבִיךָ וְעֶרְוַת אִמְּךָ לֹא תְגַלֵּה אִמְּךָ הִוא לֹא תְגַלֶּה עֶרְוָתָהּ׃}
}
{
	\eP{
		You shall not expose your father’s nudity and your
		mother’s nudity. She is your mother. You shall not expose
		her nudity.
	}
}
{
}

\Verse{8}
{
	\hP{עֶרְוַת אֵשֶׁת אָבִיךָ לֹא תְגַלֵּה עֶרְוַת אָבִיךָ הִוא׃}
}
{
	\eP{
		You shall not expose your father’s wife’s nudity. It is
		your father’s nudity.
	}
}
{
}

\Verse{9}
{
	\hP{עֶרְוַת אֲחוֹתְךָ בַת אָבִיךָ אוֹ בַת אִמֶּךָ מוֹלֶדֶת בַּיִת אוֹ מוֹלֶדֶת חוּץ לֹא תְגַלֶּה עֶרְוָתָן׃}
}
{
	\eP{
		Your sister’s nudity—your father’s daughter or your
		mother’s daughter, born home or born outside—you shall not
		expose their nudity.
	}
}
{
}

\Verse{10}
{
	\hP{עֶרְוַת בַּת בִּנְךָ אוֹ בַת בִּתְּךָ לֹא תְגַלֶּה עֶרְוָתָן כִּי עֶרְוָתְךָ הֵנָּה׃}
}
{
	\eP{
		The nudity of your son’s daughter or your daughter’s
		daughter—you shall not expose their nudity, because they
		are your nudity.
	}
}
{
}

\Verse{11}
{
	\hP{עֶרְוַת בַּת אֵשֶׁת אָבִיךָ מוֹלֶדֶת אָבִיךָ אֲחוֹתְךָ הִוא לֹא תְגַלֶּה עֶרְוָתָהּ׃}
}
{
	\eP{
		The nudity of your father’s wife’s daughter, born of your
		father: she is your sister; you shall not expose her
		nudity.
	}
}
{
}

\Verse{12}
{
	\hP{עֶרְוַת אֲחוֹת אָבִיךָ לֹא תְגַלֵּה שְׁאֵר אָבִיךָ הִוא׃}
}
{
	\eP{
		You shall not expose your father’s sister’s nudity. She is
		your father’s close relative.
	}
}
{
}

\Verse{13}
{
	\hP{עֶרְוַת אֲחוֹת אִמְּךָ לֹא תְגַלֵּה כִּי שְׁאֵר אִמְּךָ הִוא׃}
}
{
	\eP{
		You shall not expose your mother’s sister’s nudity. She is
		your mother’s close relative.
	}
}
{
}

\Verse{14}
{
	\hP{עֶרְוַת אֲחִי אָבִיךָ לֹא תְגַלֵּה אֶל אִשְׁתּוֹ לֹא תִקְרָב דֹּדָתְךָ הִוא׃}
}
{
	\eP{
		You shall not expose your father’s brother’s nudity: you
		shall not come close to his wife. She is your aunt.
	}
}
{
}

\Verse{15}
{
	\hP{עֶרְוַת כַּלָּתְךָ לֹא תְגַלֵּה אֵשֶׁת בִּנְךָ הִוא לֹא תְגַלֶּה עֶרְוָתָהּ׃}
}
{
	\eP{
		You shall not expose your daughter-in-law’s nudity. She is
		your son’s wife. You shall not expose her nudity.
	}
}
{
%	\footnote{
%		You shall not expose her nudity. The command occurs twice in the case of
%		sex with one’s daughter-in-law. It is a particularly emphatic order to
%		one who may well be in a position of sufficient power in a family to
%		commit this act. And it is a particularly emphatic prohibition of an act
%		that may be especially hurtful and psychologically damaging.
%	}
}

\Verse{16}
{
	\hP{עֶרְוַת אֵשֶׁת אָחִיךָ לֹא תְגַלֵּה עֶרְוַת אָחִיךָ הִוא׃}
}
{
	\eP{
		You shall not expose your brother’s wife’s nudity. It is
		your brother’s nudity.
	}
}
{
}

\Verse{17}
{
	\hP{עֶרְוַת אִשָּׁה וּבִתָּהּ לֹא תְגַלֵּה אֶת בַּת בְּנָהּ וְאֶת בַּת בִּתָּהּ לֹא תִקַּח לְגַלּוֹת עֶרְוָתָהּ שַׁאֲרָה הֵנָּה זִמָּה הִוא׃}
}
{
	\eP{
		You shall not expose the nudity of a woman and her
		daughter. You shall not take her son’s daughter or her
		daughter’s daughter to expose her nudity. They are close
		relatives. It is perversion.
	}
}
{
}

\Verse{18}
{
	\hP{וְאִשָּׁה אֶל אֲחֹתָהּ לֹא תִקָּח לִצְרֹר לְגַלּוֹת עֶרְוָתָהּ עָלֶיהָ בְּחַיֶּיהָ׃}
}
{
	\eP{
		And you shall not take a woman to her sister to rival, to
		expose her nudity along with her in her lifetime.
	}
}
{
%	\footnote{
%		along with her. The Hebrew and this English translation both contain an
%		ambiguity: does it mean to marry two sisters (as it is frequently
%		understood), to have sexual relations (whether married or not) with two
%		sisters, or to have sexual relations with two sisters at the same time
%		(as Rashi understands it)? It says, “in her lifetime”—which sounds as
%		though it means marriage.
%	}
}

\Verse{19}
{
	\hP{וְאֶל אִשָּׁה בְּנִדַּת טֻמְאָתָהּ לֹא תִקְרַב לְגַלּוֹת עֶרְוָתָהּ׃}
}
{
	\eP{
		And you shall not come close to a woman during her
		menstrual impurity to expose her nudity.
	}
}
{
}

\Verse{20}
{
	\hP{וְאֶל אֵשֶׁת עֲמִיתְךָ לֹא תִתֵּן שְׁכבְתְּךָ לְזָרַע לְטמְאָה בָהּ׃}
}
{
	\eP{
		And you shall not give your intercourse of seed to your
		fellow’s wife, to become impure by it.
	}
}
{
%	\footnote{
%		give your intercourse of seed. Why is the wording different in this
%		verse from those that precede it—i.e., “give your intercourse of seed”
%		instead of “expose nudity”? Perhaps it is to connect to the next verse.
%		See the next comment.
%	}
}

\Verse{21}
{
	\hP{וּמִזַּרְעֲךָ לֹא תִתֵּן לְהַעֲבִיר לַמֹּלֶךְ וְלֹא תְחַלֵּל אֶת שֵׁם אֱלֹהֶיךָ אֲנִי יְהֹוָה׃}
}
{
	\eP{
		And you shall not give any of your seed for passing to the
		Molech, and you shall not desecrate your God’s name. I am
		YHWH.
	}
}
{
%	\footnote{
%		give any of your seed. It has been asked why this commandment is placed
%		here, grouped with the commandments against particular sexual unions. It
%		must be relevant that this commandment and the commandment against
%		adultery that precedes it both use the expression “to give seed.”
%		Perhaps the command against adultery and the command against sacrificing
%		one’s children came to be connected precisely because of the parallel of
%		the wording in which each had come to be known. Or perhaps it is the
%		reverse: that the “seed” wording was used in the preceding verse
%		precisely to make the transition to the “Molech” law. A third
%		possibility is that the point is precisely the seed: that sexual
%		intercourse with a married woman involves a potential abuse of children,
%		and the fact that the law against child sacrifice is adjacent to this
%		prohibition serves to emphasize this danger. Footnote 18:21. for passing
%		to the Molech. This “passing” must be understood in the light of the
%		earlier commandment that requires the sacrifice of every firstborn: “you
%		will pass every first birth of a womb to YHWH” (Exod 13:12). There, the
%		Israelite is told to explain to his children: “I am sacrificing to YHWH
%		every first birth of a womb, and I shall redeem every firstborn of my
%		sons” (13:15). We learn from this that the term “passing” in such
%		contexts refers to sacrifice and that Israelites are not permitted to
%		sacrifice their children, either to the God of Israel or to a foreign
%		deity. This is supported by a reference in 2 Kings 23:10 to “a man’s
%		passing his son or his daughter in fire to the Molech.” Footnote 18:21.
%		Molech. Some have taken this to be the name of a pagan god to whom human
%		sacrifices were made. But there is little evidence of any human
%		sacrifice in the ancient Near East, and no evidence of such sacrifice to
%		a god named Molech (or any similar name). There is evidence of human
%		sacrifice at the Phoenician colony at Carthage, but there it is to the
%		god Baal-Hammon. There is evidence that the term Molech refers rather to
%		the institution, the practice of human sacrifice. The law here would
%		thus mean that one should not bring one’s offspring for the practice of
%		Molech. The reference in Lev 20:2 below may shed further light on this.
%	}
}

\Verse{22}
{
	\hP{וְאֶת זָכָר לֹא תִשְׁכַּב מִשְׁכְּבֵי אִשָּׁה תּוֹעֵבָה הִוא׃}
}
{
	\eP{
		And you shall not lie with a male like lying with a woman.
		It is an offensive thing.
	}
}
{
%	\footnote{
%		you shall not lie with a male like lying with a woman. Why is male
%		homosexuality explicitly forbidden in the Torah but not female? Some
%		would surmise that it is because women are controlled in a patriarchal
%		Israelite society; and so a woman would simply have no choice but to
%		marry a man. But this is not an adequate explanation, because there
%		would still be opportunities for female homosexual liaisons. Some would
%		say that the concern is the seed, which is understood to come from the
%		male, and therefore is “wasted” in another male. But the text calls
%		homosexuality “an offensive thing” (in older translations: “an
%		abomination”), which certainly sounds like an abhorrence of the act, and
%		not just a concern with the practical matter of reproduction. The reason
%		may rather be because the Torah comes from a world in which there is
%		polygamy. A man can have sex with his two wives simultaneously. That
%		this is understood to be permissible is implied by the fact that the law
%		in v. 18 above forbids it only with sisters (see the comment). Or, even
%		if the above case means marriage and not simultaneous sex, then
%		simultaneous sex still is not forbidden anywhere in the Torah. If
%		simultaneous sex with one’s two (or more) wives is practiced, it would
%		be difficult to allow this while forbidding female homosexuality. (At
%		minimum, it could require a number of laws specifying what sort of
%		contact is permissible and under what circumstances.) In the present
%		state of knowledge concerning homosexuality, it is difficult to justify
%		its prohibition in the Torah. All of the movements in Judaism (and other
%		religions) are currently contending with this issue. Its resolution
%		ultimately must lie in the law of Deuteronomy that states that, for
%		difficult matters of the law, people must turn to the authorities of
%		their age, to those who are competent to judge, and those judges must
%		decide (Deut 17:8–9). In my own view, the present understanding of the
%		nature of homosexuality indicates that it is not an “offensive thing”
%		(also translated “abomination”) as described in this verse. The Hebrew
%		term for ”offensive thing” (tô‘bh) is understood to be a relative term,
%		which varies according to human perceptions. For example, in Genesis,
%		Joseph tells his brothers that “any shepherd is an offensive thing to
%		Egypt” (46:34); but, obviously, it is not an offensive thing to the
%		Israelites. In light of the evidence at present, homosexuality cannot be
%		said to be unnatural, nor is it an illness. Its prohibition in this
%		verse explicitly applies only so long as it is properly perceived to be
%		offensive, and therefore the current state of the evidence suggests that
%		the period in which this commandment was binding has come to an end.
%	}
}

\Verse{23}
{
	\hP{וּבְכל בְּהֵמָה לֹא תִתֵּן שְׁכבְתְּךָ לְטמְאָה בָהּ וְאִשָּׁה לֹא תַעֲמֹד לִפְנֵי בְהֵמָה לְרִבְעָהּ תֶּבֶל הוּא׃}
}
{
	\eP{
		And you shall not give your intercourse with any animal,
		to become impure by it. And a woman shall not stand in
		front of an animal to mate with it. It is an aberration.
	}
}
{
}

\Verse{24}
{
	\hP{אַל תִּטַּמְּאוּ בְּכל אֵלֶּה כִּי בְכל אֵלֶּה נִטְמְאוּ הַגּוֹיִם אֲשֶׁר אֲנִי מְשַׁלֵּחַ מִפְּנֵיכֶם׃}
}
{
	\eP{
		“You shall not become impure by all of these, because the
		nations that I am putting out from in front of you became
		impure by all of these,
	}
}
{
}

\Verse{25}
{
	\hP{וַתִּטְמָא הָאָרֶץ וָאֶפְקֹד עֲוֺנָהּ עָלֶיהָ וַתָּקִא הָאָרֶץ אֶת יֹשְׁבֶיהָ׃}
}
{
	\eP{
		and the land became impure, and I reckoned its sin on it,
		and the land vomited out its residents.
	}
}
{
}

\Verse{26}
{
	\hP{וּשְׁמַרְתֶּם אַתֶּם אֶת חֻקֹּתַי וְאֶת מִשְׁפָּטַי וְלֹא תַעֲשׂוּ מִכֹּל הַתּוֹעֵבֹת הָאֵלֶּה הָאֶזְרָח וְהַגֵּר הַגָּר בְּתוֹכְכֶם׃}
}
{
	\eP{
		But you: you shall observe my laws and my judgments and
		not do any of all these offensive things, the citizen and
		the alien who resides among you—
	}
}
{
}

\Verse{27}
{
	\hP{כִּי אֶת כּל הַתּוֹעֵבֹת הָאֵל עָשׂוּ אַנְשֵׁי הָאָרֶץ אֲשֶׁר לִפְנֵיכֶם וַתִּטְמָא הָאָרֶץ׃}
}
{
	\eP{
		because the people of the land who were before you did all
		of these offensive things, and the land became impure—
	}
}
{
}

\Verse{28}
{
	\hP{וְלֹא תָקִיא הָאָרֶץ אֶתְכֶם בְּטַמַּאֲכֶם אֹתָהּ כַּאֲשֶׁר קָאָה אֶת הַגּוֹי אֲשֶׁר לִפְנֵיכֶם׃}
}
{
	\eP{
		so the land will not vomit you out for your making it
		impure as it vomited out the nation that was before you.
	}
}
{
}

\Verse{29}
{
	\hP{כִּי כּל אֲשֶׁר יַעֲשֶׂה מִכֹּל הַתּוֹעֵבֹת הָאֵלֶּה וְנִכְרְתוּ הַנְּפָשׁוֹת הָעֹשֹׂת מִקֶּרֶב עַמָּם׃}
}
{
	\eP{
		Because anyone who will do any of these offensive things:
		the persons who do will be cut off from among their
		people.
	}
}
{
}

\Verse{30}
{
	\hP{וּשְׁמַרְתֶּם אֶת מִשְׁמַרְתִּי לְבִלְתִּי עֲשׂוֹת מֵחֻקּוֹת הַתּוֹעֵבֹת אֲשֶׁר נַעֲשׂוּ לִפְנֵיכֶם וְלֹא תִטַּמְּאוּ בָּהֶם אֲנִי יְהֹוָה אֱלֹהֵיכֶם׃}
}
{
	\eP{
		And you shall keep my charge: not to do any of the
		abominable customs that were done before you; and you will
		not become impure by them. I am YHWH, your God.”
	}
}
{
}
